% Part: proof-theory
% Chapter: natural-deduction
% Section: quantifiers

\documentclass[../../../include/open-logic-section]{subfiles}

\begin{document}

\olfileid{pt}{nat}{qua}

\olsection{નિયમિત \usetoken{P}{proof} અને પ્રતિસ્થાપન}

પરિમાણકોના નિયમોમાં
\Elim{\lexists} અને
\Intro{\lforall}ને લાગતી
``સ્વનિર્ધારક ચલની શરતો''
ને કારણે કેટલીક મુશ્કેલીઓ
ઊભી થાય છે. આ અનુમાનોમાં
આવતો !!{constant}~$c$
ફરી કદી ન વપરાય એ
સુનિશ્ચિત કરીએ તો
મોટાભાગની મુશ્કેલીઓ
ઉકેલી શકાય છે.
આ માટે નીચેની વ્યાખ્યા લો.

\begin{defn}
સ્વાભાવિક નિગમનની
!!{proof}~$\delta$
\emph{નિયમિત} છે, જો
\begin{enumerate}
  \item કોઈ સ્વનિર્ધારક ચલ~$a$
  એકથી વધુ \Elim{\lexists}
  અથવા~\Intro{\lforall}
  અનુમાનમાં સ્વનિર્ધારક ચલ
  તરીકે ન વપરાતો હોય, અને
  \item જો~$a$ કોઈ
  \Elim{\lexists} અથવા
  \Intro{\lforall} અનુમાનનો
  સ્વનિર્ધારક ચલ હોય,
  તો તે અનુક્રમે ગૌણ
  આધારવિધાન અથવા આધારવિધાન
  સુધી પહોંચતી તાત્કાલિક
  ઉપ-!!{proof}માં જ આવતો હોય.
\end{enumerate}
\end{defn}

\begin{lem}\ollabel{lem:subst-regular}
જો~$\delta$ નિયમિત હોય,
તેનો અંત~$!C$માં થતો હોય,
$c$ એવો !!a{constant} હોય
જે~$\delta$માં સ્વનિર્ધારક
ચલ તરીકે વપરાયો ન હોય,
અને~$t$ એવું પદ હોય
જેમાં~$\delta$નો કોઈ
સ્વનિર્ધારક ચલ ન હોય,
તો~$\delta$માં~$c$ના
દરેક સ્થાનને~$t$થી
બદલવાથી મળતી
$\Subst{\delta}{t}{c}$
નિયમિત !!{proof} છે.
$\Subst{\delta}{t}{c}$નું
અંતિમ-!!{formula}
$\Subst{!C}{t}{c}$ છે.
જો~$!A^x$ એ~$\delta$ની
ખુલ્લી ધારણા હોય, તો
$\Subst{!A}{t}{c}^x$ એ
$\Subst{\delta}{t}{c}$ની
ખુલ્લી ધારણા છે.
\end{lem}

\begin{proof}
  $\pheight{\delta}$ પર
  આગમન કરીએ.
  જો~$\pheight{\delta} = 0$,
  તો તેમાં માત્ર
  ધારણા~$!A^x$ છે.
  ત્યારે~$\Subst{\delta}{t}{c}$
  એ~$\Subst{!A}{t}{c}^x$ છે.

જો~$\pheight{\delta} > 0$,
  તો~$\delta$ના છેલ્લા
  અનુમાન પ્રમાણે કિસ્સા
  જુદા પાડીએ.
  વિધાનાત્મક સંયોજકોના
  નિયમોના કિસ્સા સહેલા છે
  અને કસરત તરીકે રાખ્યા છે.
  અહીં~$\delta$નો અંત
  પરિમાણક-અનુમાનમાં થાય
  તે કિસ્સા જોઈએ.
  \begin{enumerate}
    \item છેલ્લું અનુમાન
    $\Intro{\lforall}$ હોય,
    તો~$\delta$નું સ્વરૂપ
    \[
    \AxiomC{}
    \RightLabel{$\delta'$}
    \DeduceC{$!A(a,c)$}
    \RightLabel{\Intro{\lforall}}
    \UnaryInfC{$\lforall[x][!A(x,c)]$}
    \DisplayProof
    \]
    આગમન-પરિકલ્પનાથી
    $\Subst{\delta'}{t}{c}$
    એ~$!A(a,t)$ની નિયમિત
    !!{proof} છે.
    $a$ એ~$\delta$નો
    સ્વનિર્ધારક ચલ હોવાથી
    $t$માં~$a$ આવતો નથી.
    તેથી~$\Subst{\delta}{t}{c}$નું
    છેલ્લું અનુમાન,
    \[
    \AxiomC{}
    \RightLabel{$\Subst{\delta'}{t}{c}$}
    \DeduceC{$!A(a,t)$}
    \RightLabel{\Intro{\lforall}}
    \UnaryInfC{$\lforall[x][!A(x,t)]$}
    \DisplayProof
    \]
    યોગ્ય છે: $t$માં~$a$
    ન હોવાથી નિષ્કર્ષમાં
    કે કોઈ ખુલ્લી ધારણામાં
    $a$ આવતો નથી.
    \item છેલ્લું અનુમાન
    $\Elim{\lforall}$ હોય,
    તો~$\delta$નું સ્વરૂપ
    \[
    \AxiomC{}
    \RightLabel{$\delta'$}
    \DeduceC{$\lforall[x][!A(x,c)]$}
    \RightLabel{\Elim{\lforall}}
    \UnaryInfC{$!A(s(c),c)$}
    \DisplayProof
    \]
    આગમન-પરિકલ્પનાથી
    $\Subst{\delta'}{t}{c}$
    એ~$\lforall[x][!A(x,t)]$ની
    નિયમિત !!{proof} છે.
    (નિષ્કર્ષમાંનું
    પદ~$s(c)$, $c$ ધરાવતું
    હોય પણ શકે અને ન પણ હોય.)
    $\Subst{\delta}{t}{c}$નું
    છેલ્લું અનુમાન,
    \[
    \AxiomC{}
    \RightLabel{$\Subst{\delta'}{t}{c}$}
    \DeduceC{$\lforall[x][!A(x,t)]$}
    \RightLabel{\Elim{\lforall}}
    \UnaryInfC{$!A(s(t),t)$}
    \DisplayProof
    \]
    યોગ્ય છે, અને
    $!A(s(t),t) \ident \Subst{!A(s(c),c)}{t}{c}$.
    \Intro{\lexists}નો
    કિસ્સો આને સમાન છે.
  \item છેલ્લું અનુમાન
    $\Elim{\lexists}$ હોય,
    તો~$\delta$નું સ્વરૂપ
    \[
    \AxiomC{}
    \RightLabel{$\delta_1'$}
    \DeduceC{$\lexists[x][!A(x,c)]$}
    \AxiomC{$\Discharge{!A(a,c)}{x}$}
    \RightLabel{$\delta_2'$}
    \DeduceC{$!C$}
    \DischargeRule{\Elim{\lexists}}{x}
    \BinaryInfC{$!C$}
    \DisplayProof
    \]
    આગમન-પરિકલ્પનાથી
    $\Subst{\delta_1'}{t}{c}$
    અને~$\Subst{\delta_2'}{t}{c}$
    અનુક્રમે
    $\lexists[x][!A(x,t)]$ની
    અને ખુલ્લી ધારણા
    $!A(a,t)$માંથી
    $\Subst{!C}{t}{c}$ની
    નિયમિત !!{proof}s છે.
    $a$ એ~$\delta$નો
    સ્વનિર્ધારક ચલ હોવાથી
    $t$માં~$a$ આવતો નથી.
    આથી~$\Subst{\delta}{t}{c}$નું
    છેલ્લું અનુમાન,
    \[
    \AxiomC{}
    \RightLabel{$\Subst{\delta_1'}{t}{c}$}
    \DeduceC{$\lexists[x][!A(x,t)]$}
    \AxiomC{$\Discharge{!A(a,t)}{x}$}
    \RightLabel{$\Subst{\delta_2'}{t}{c}$}
    \DeduceC{$\Subst{!C}{t}{c}$}
    \DischargeRule{\Elim{\lexists}}{x}
    \BinaryInfC{$\Subst{!C}{t}{c}$}
    \DisplayProof
    \]
    યોગ્ય છે:
    $\Subst{!A(x,t)}{a}{x} \ident \Subst{!A(a,c)}{t}{c}$,
    અને નિષ્કર્ષ
    $\Subst{!C}{t}{c}$માં
    કે કોઈ ખુલ્લી ધારણામાં
    $a$ આવતો નથી.
    \sourcecorrection{OLMD-228}{મૂળના અસ્તિત્વ-નિવારણ કિસ્સામાં પ્રતિસ્થાપિત પહેલા નિષ્કર્ષમાં $c$ જ રહી ગયું છે, બીજા નિષ્કર્ષમાં $C$નું પ્રતિસ્થાપન છૂટ્યું છે અને અંતિમ વૃક્ષ ભૂલથી સર્વમાન્યક-પરિચયનું છે; અહીં બંને પ્રતિસ્થાપિત ઉપસાબિતીથી અસ્તિત્વ-નિવારણનું યોગ્ય વૃક્ષ મૂક્યું છે.}
  \end{enumerate}
\end{proof}

\begin{prop}\ollabel{prop:regular}
જો~$\delta$ એ~\Log{N1c}
અથવા~\Log{N1i}માંની
!!a{proof} હોય,
તો એ જેવી જ
રચના ધરાવતી
(ખાસ કરીને સરખી
!!{height} ધરાવતી)
નિયમિત !!{proof}~$\delta'$
છે, જે~$\delta$થી
માત્ર સ્વનિર્ધારક ચલો
બદલવામાં જુદી પડે છે.
\end{prop}

\begin{proof}
આ સાબિતી માટે~$\delta$માંના
સ્વનિર્ધારક ચલ ધરાવતા
અનુમાનને \emph{સ્વચ્છ}
કહીએ, જો તેનો સ્વનિર્ધારક
ચલ બીજા કોઈ અનુમાનનો
સ્વનિર્ધારક ચલ ન હોય
અને તે આ અનુમાનની
તાત્કાલિક ઉપ-!!{proof}
સિવાય~$\delta$માં ક્યાંય
ન આવતો હોય; નહિતર તેને
\emph{અસ્વચ્છ} કહીએ.
$n$ એ~$\delta$માંનાં
અસ્વચ્છ અનુમાનોની સંખ્યા
હોય. $n$ પર આગમનથી
બતાવીશું કે~$\delta$ને
જરૂરી નિયમિત
!!{proof}~$\delta'$માં
ફેરવી શકાય છે.
જો~$n=0$, તો~$\delta$નાં બધાં
સ્વનિર્ધારક ચલવાળાં અનુમાનો
સ્વચ્છ છે, એટલે~$\delta$
નિયમિત છે અને કંઈ સાબિત
કરવાનું રહેતું નથી.
નહિતર સૌથી ઉપરનું
અસ્વચ્છ સ્વનિર્ધારક ચલવાળું
અનુમાન પસંદ કરો;
માનો કે તે~\Intro{\lforall}
છે. તેનામાં પૂરી થતી
$\delta$ની
ઉપ-!!{proof}~$\delta_1$
આ સ્વરૂપની છે:
    \[
    \AxiomC{}
    \RightLabel{$\delta_2$}
    \DeduceC{$!A(c)$}
    \RightLabel{\Intro{\lforall}}
    \UnaryInfC{$\lforall[x][!A(x)]$}
    \DisplayProof
    \]
$c$ સ્વનિર્ધારક ચલ હોવાથી
તે~$\lforall[x][!A(x)]$માં
કે કોઈ ખુલ્લી ધારણામાં
આવતો નથી.
તેમાં અસ્વચ્છ
સ્વનિર્ધારક ચલવાળું અનુમાન
નથી, તેથી~$\delta_2$
નિયમિત છે. વળી પસંદ કરેલો
ચલ તેના કોઈ અનુમાનનો
સ્વનિર્ધારક ચલ નથી;
હોય તો તે અનુમાન પણ
અસ્વચ્છ બનત.
\sourcecorrection{OLMD-229}{મૂળમાં સૌથી ઉપરનું કોઈ પણ સ્વનિર્ધારક-ચલવાળું અનુમાન પસંદ થાય છે અને ઉપરની ઉપસાબિતીમાં આવાં અનુમાનો જ નથી એમ કહેવાયું છે; આગમન માટે સૌથી ઉપરનું \emph{અસ્વચ્છ} અનુમાન જોઈએ, જેથી તેની ઉપર સ્વચ્છ અનુમાનો રહી શકે પણ ઉપસાબિતી નિયમિત રહે.}
$\delta$માં ન આવતો
!!a{constant}~$a$ લો.
\olref{lem:subst-regular}થી
$\Subst{\delta_2}{a}{c}$
એ~$!A(a)$ની નિયમિત
સાબિતી છે.
સ્વનિર્ધારક ચલની શરતથી
$\delta_2$ની કોઈ ખુલ્લી
ધારણામાં~$c$ નથી,
તેથી~$\Subst{\delta_2}{a}{c}$ની
કોઈ ખુલ્લી ધારણામાં~$a$
નથી. આથી~$\Intro{\lforall}$
લાગુ કરી શકાય છે.
જો~$\delta$માં
$\delta_1$ને બદલે
સાબિતી~$\delta_1'$ મૂકીએ,
    \[
    \AxiomC{}
    \RightLabel{$\Subst{\delta_2}{a}{c}$}
    \DeduceC{$!A(a)$}
    \RightLabel{\Intro{\lforall}}
    \UnaryInfC{$\lforall[x][!A(x)]$}
    \DisplayProof
    \]
તો અસ્વચ્છ સ્વનિર્ધારક
ચલવાળું અનુમાન સ્વચ્છ
અનુમાનથી બદલાયું છે,
અને ફરક માત્ર
સ્વનિર્ધારક ચલ~$c$ની
જગ્યાએ~$a$ મૂકવાનો છે.
મળેલી સાબિતીમાં
વધુમાં વધુ~$n-1$
અસ્વચ્છ સ્વનિર્ધારક
ચલવાળાં અનુમાનો છે;
આગમન-પરિકલ્પનાથી
પરિણામ મળે છે.
\sourcecorrection{OLMD-230}{મૂળમાં અસ્વચ્છ અનુમાનોની સંખ્યા ચોક્કસ $n-1$ થાય છે એમ કહેવાયું છે; તાજા ચલથી બીજાં અનુમાનો પણ સ્વચ્છ થઈ શકે, તેથી આગમન માટે જરૂરી વધુમાં વધુ $n-1$ દાવો લખ્યો છે.}
\end{proof}

\begin{prob}
\olref{prop:regular}ની
સાબિતીમાં સૌથી ઉપરનું
અસ્વચ્છ સ્વનિર્ધારક
ચલવાળું અનુમાન~\Elim{\lexists}
હોય તે કિસ્સો વર્ણવો.
\end{prob}

હમણાં સાબિત કરેલી
પ્રતિજ્ઞાપનનો સ્પષ્ટ
ઉલ્લેખ હવે નહીં કરીએ;
પણ ચર્ચામાં આવતી
બધી !!{proof}s નિયમિત
છે એમ માનીશું.
નિયમિત ન હોય તો
સ્વનિર્ધારક ચલો બદલીને
તેમને નિયમિત બનાવી
શકાય છે.

\Olref{lem:subst-regular}
અને~\olref{prop:regular}
\Log{N2c} તથા~\Log{N2i}
માટે પણ સાચાં છે.
તેમની સાબિતીઓ કસરત
તરીકે રાખીએ છીએ.

\end{document}
